% NexusQuant — Introduction
% NeurIPS 2026 submission
% ~1.5 pages. One clear story: eviction + E8 VQ, training-free, GPU-validated.
% All comparison numbers cite published papers; no running competitors' code.

%------------------------------------------------------------------
\section{Introduction}
\label{sec:intro}
%------------------------------------------------------------------

Context windows are growing. GPT-4o handles 128K tokens; Gemini 1.5 reaches
1M. Yet memory grows with them: at 128K tokens, a 70B-parameter model
accumulates over 40\,GB in its key-value (KV) cache alone---exceeding the
capacity of a single A100 and making multi-turn long-context inference
prohibitively expensive for most deployments. Compressing the KV cache is
therefore one of the highest-impact open problems in LLM systems research.

\paragraph{The bottleneck.}
Formally, the KV cache for a model with $L$ layers, $H$ KV heads, head
dimension $d$, and sequence length $T$ occupies $2 \cdot L \cdot H \cdot d
\cdot T$ floating-point values. For a 7B model in FP16, this grows at
$\approx$\,0.5\,MB per token, making 128K contexts cost $\approx$\,64\,GB.
Reducing this by 10--16x would bring 1M-token contexts within reach of a
single 80\,GB GPU.

\paragraph{Two complementary redundancies.}
Prior work has attacked this problem along two orthogonal axes, rarely
combined:

\emph{Token eviction} (H2O~\cite{h2o2023}, SnapKV~\cite{snapkv2024}) removes
the KV vectors of tokens deemed unimportant by attention scores. At aggressive
rates, eviction alone causes catastrophic quality loss because important
information is discarded permanently.

\emph{Quantization} (KIVI~\cite{kivi2024}, KVQuant~\cite{kvquant2024},
TurboQuant~\cite{turboquant2026}) reduces the precision of KV vectors while
retaining all tokens. The best training-free method, TurboQuant (Google,
ICLR 2026), achieves $\approx$\,5--6x compression at near-zero quality loss
using polar-coordinate quantization and Johnson--Lindenstrauss projection.
To exceed $\sim$\,6x training-free, scalar and vector quantization alone
run out of room.

\paragraph{Our insight: combine eviction and quantization.}
Eviction reduces the \emph{number} of tokens retained; quantization reduces
the \emph{precision} of each retained vector. These redundancies are nearly
orthogonal: a 50\% eviction halves token count independently of how many bits
are used per token. Combining 60\% eviction with 2-bit E8 lattice vector
quantization yields 16.6x compression---a regime no prior training-free
method reaches.

The key challenge is that naive importance-based eviction interacts badly with
quantization noise: short contexts ($\le$\,500 tokens) provide too few tokens
for the attention scorer to reliably distinguish important from unimportant
entries, and the eviction cliff (the rate at which quality degrades as more
tokens are dropped) is steep. We find that this cliff shifts dramatically with
context length.

\paragraph{Context-length scaling: a novel phenomenon.}
Our central empirical discovery is that \emph{quality improves with longer
context}, contrary to the common assumption that longer sequences are harder to
compress. At 500-token prefixes, 35\% eviction gives $+0.90\%$ PPL degradation
and the cliff appears at 36\% eviction. At 2048-token prefixes on the same
model, 35\% eviction gives only $+0.14\%$ PPL, 60\% eviction gives $+0.82\%$
($<$\,1\%), and even 80\% eviction stays at $+2.1\%$---a 5--6x improvement in
quality at every eviction rate. This scaling occurs because longer prefixes
give the attention-based importance scorer more signal: with hundreds of query
positions to aggregate over, the scorer learns which tokens matter and drops
the right ones. This finding is practically significant: production inference
typically operates at thousands of tokens, where our compression ratios are
substantially better than short-context benchmarks suggest.

\paragraph{Contributions.}
We make the following contributions:

\begin{enumerate}
  \item \textbf{NexusQuant pipeline.} A training-free, zero-calibration
  compression pipeline: attention importance scoring $\to$ token eviction
  $\to$ RoPE removal from keys $\to$ Hadamard incoherence preprocessing $\to$
  2-bit E8 lattice vector quantization $\to$ temporal delta coding $\to$ zstd
  entropy compression. All compression ratios are measured via
  \texttt{torch.cuda.memory\_allocated()} with all overhead included.

  \item \textbf{Context-length scaling law.} We document and quantify the
  phenomenon that attention-aware eviction quality scales with prefix length,
  improving by 5--6x between 500 and 2048 tokens. This explains why
  short-context PPL benchmarks understate the practical utility of
  eviction-based methods.

  \item \textbf{Training-free frontier at 16.6x.} On Mistral-7B with 2048-token
  prefixes, we achieve 10.4x at $+0.14\%$ PPL, 13.4x at $+0.56\%$, and
  \textbf{16.6x at $+0.82\%$}---the first training-free result to exceed 16x
  at $<$\,1\% PPL on a 7B-class model without calibration data.

  \item \textbf{E8 lattice entropy advantage.} E8 lattice indices compress 3x
  better than scalar quantization indices under entropy coding (zstd level 22),
  because the E8 nearest-neighbor assignment clusters indices non-uniformly,
  creating a highly skewed distribution amenable to Huffman/ANS coding.

  \item \textbf{Honest domain characterization.} We evaluate on three text
  domains (academic, technical, creative/narrative) and find that creative text
  is the hardest case at low eviction rates ($+2.5\%$ vs. $+0.14\%$ for
  academic at 35\% eviction, long context), with convergence at high eviction
  rates ($\approx$\,4--5\% across all domains at 80\% eviction). This domain
  sensitivity should be reported by all future eviction-based methods.
\end{enumerate}

\paragraph{What we do not claim.}
We compare against published numbers for KVTC, TurboQuant, CommVQ, and Palu;
we do not re-run their code on our evaluation data, so those comparisons are
directional rather than controlled. We evaluate on Mistral-7B and TinyLlama;
Llama-3 results are pending. Our evaluation metric is perplexity on
WikiText-103 and multi-topic text; we do not report LongBench task accuracy in
this version. Compression ratios are measured from KV cache memory only; we do
not report end-to-end inference latency.

\paragraph{Paper organization.}
Section~\ref{sec:related} reviews related work.
Section~\ref{sec:method} describes the pipeline in detail.
Section~\ref{sec:experiments} presents GPU-validated results on Mistral-7B,
including the context-length scaling analysis, domain sensitivity, and
ablation study.
Section~\ref{sec:limitations} discusses limitations and future work.
